\section{问题重述}

\subsection{问题背景}

某医院超声科同时接收住院、门诊和体检三类检查申请。一次申请可能包含多个项目，各项目在检查时长、准备条件和可用设备上存在差异；22台设备具有明确的专项分工，医生又在不同设备之间共享。任一时刻，一台设备只能承担一项检查，同一患者的多个项目也不能并行执行。住院患者还受到开单后48小时完成全部项目的时限约束，门诊与体检需求则集中占用计划日内的服务容量。医疗预约研究通常需要同时处理等待、容量和服务水平\cite{gupta2008,ahmadi2017}，分时预约和检查资源协同也是改善候检秩序的重要手段\cite{nhc2020,chenjuan2023}，因此需要在设备能力、人员容量与患者完整流程之间进行统一配置。

\subsection{问题重述}

根据题目给出的历史检查记录、设备项目说明和医生信息，需要分析超声需求与资源状态，并针对纯住院场景及三类患者共享资源场景制定预约优化方案。具体需要完成以下三个问题：

（1）分析住院、门诊和体检需求随日期、科室与项目的变化，识别平峰、高峰及压力状态，估计项目工时和医生并发能力，并据此给出设备开放、专项资源保留和预约分流原则。

（2）仅考虑住院患者，在48小时内完整完成全部检查的口径下，确定患者主序、项目次序、检查时刻和所用设备，使完成事件数尽可能多，并输出患者级预约与检查顺序。

（3）在问题二的资源约束基础上加入门诊和体检患者，建立三类患者联合调度模型，刻画背景服务率与住院48小时完成率之间的交换关系，选择可实施的折中方案并检验其稳定性。
